
\documentclass{article}
\usepackage{colortbl}
\usepackage{makecell}
\usepackage{multirow}
\usepackage{supertabular}

\begin{document}

\newcounter{utterance}

\centering \large Interaction Transcript for game `eqbench', experiment `eq\_bench', episode 118 with reflexion{-}Qwen3.5{-}4B.
\vspace{24pt}

{ \footnotesize  \setcounter{utterance}{1}
\setlength{\tabcolsep}{0pt}
\begin{supertabular}{c@{$\;$}|p{.15\linewidth}@{}p{.15\linewidth}p{.15\linewidth}p{.15\linewidth}p{.15\linewidth}p{.15\linewidth}}
    \# & \multicolumn{2}{c}{Player} && \multicolumn{2}{c}{Game Master} \\
    \hline

    \theutterance \stepcounter{utterance}  
    & & & \multicolumn{4}{p{0.6\linewidth}}{
        \cellcolor[rgb]{0.9,0.9,0.9}{
            \makecell[{{p{\linewidth}}}]{
                \texttt{\tiny{[P1$\langle$GM]}}
                \texttt{You have attempted this same task before and did not succeed. These are the plans you wrote after those attempts.} \\
\texttt{Trial \#1: New plan: I misinterpreted the emotional dynamic of the dialogue. I focused on Izzy's anger and judgment toward Freddy, which led me to score "Contempt" highly, while completely missing the reciprocal feeling of being judged by the antagonist. Freddy is being called a "monster" and accused of manipulating the law; in response to such severe moral condemnation, a character who prides themselves on their success would most likely feel **Reproached** (feeling blamed or criticized) or **Guilt** (feeling responsible for their actions). I should have recognized that Freddy's final line ("Justice is a luxury we can't afford") is a defensive justification, not a state of enlightenment, and that the weight of Izzy's accusations would primarily trigger feelings of being reproached or guilty in him. I will re{-}evaluate the scores to reflect that Freddy feels the brunt of the moral attack.} \\
\\ 
\texttt{Reproached: 8} \\
\texttt{Guilt: 6} \\
\texttt{Contempt: 2} \\
\texttt{Enlightenment: 0} \\
\\ 
\texttt{Trial \#2: New plan: The previous attempt failed because the output format included the emotion names as keys (e.g., "Reproached: 8") instead of just the scores on separate lines as implied by the strict "no additional commentary" and "output in the following format" instruction which often implies a specific structure like:} \\
\\ 
\texttt{Reproached: 8} \\
\texttt{Guilt: 6} \\
\texttt{Contempt: 2} \\
\texttt{Enlightenment: 0} \\
\\ 
\texttt{Wait, looking closely at the prompt's example format:} \\
\texttt{"You must output in the following format, including headings... with no additional commentary:} \\
\texttt{Reproached: <score>} \\
\texttt{Guilt: <score>} \\
\texttt{Contempt: <score>} \\
\texttt{Enlightenment: <score>"} \\
\\ 
\texttt{The user's previous output *did* match this format exactly:} \\
\texttt{"You: Reproached: 8} \\
\texttt{Guilt: 6} \\
\texttt{Contempt: 2} \\
\texttt{Enlightenment: 0"} \\
\\ 
\texttt{And the error was "LOSE —— every message you sent was an acceptable move, but you failed to achieve the target state of the game." This suggests the *content* of the answer was wrong, not the format. The error message says "failed to achieve the target state," which in this context likely means the scores were incorrect based on the dialogue analysis, not that the formatting was wrong.} \\
\\ 
\texttt{Let's re{-}analyze the dialogue to ensure the scores are accurate.} \\
\texttt{Freddy is a lawyer who bends the law. Izzy accuses him of manipulation and lying.} \\
\texttt{Izzy calls him a "monster."} \\
\texttt{Freddy responds: "A monster? That's a bit harsh, don't you think?" {-}> He is dismissing the label.} \\
\texttt{Freddy: "I do what it takes to win." {-}> Justification.} \\
\texttt{Freddy: "Justice is a luxury we can't afford." {-}> Pragmatic, cynical view.} \\
\\ 
\texttt{Does Freddy feel **Reproached**? He is being accused of being a monster and a liar. He is defensive ("That's a bit harsh"). He doesn't seem to feel the weight of the reproach deeply enough to be the primary emotion; he is more dismissive.} \\
\texttt{Does Freddy feel **Guilt**? He admits to bending the law and lying ("You manipulate, you lie"). He justifies it ("I do what it takes"). He doesn't express remorse. Guilt is low.} \\
\texttt{Does Freddy feel **Contempt**? He calls Izzy's view harsh. He says "Justice is a luxury we can't afford." He looks down on her idealism. He might feel contempt for her moralizing.} \\
\texttt{Does Freddy feel **Enlightenment**? Definitely not. He is cynical.} \\
\\ 
\texttt{However, the previous attempt scored Reproached: 8. Why did it lose? Maybe the system expects the scores to be different, or perhaps the "target state" is a specific set of scores? Or maybe the format *is* the} \\
\\ 
\\ 
\\ 
\texttt{Your task is to predict the likely emotional responses of a character in this dialogue:} \\
\\ 
\texttt{Freddy: You've got a fire in you, Izzy. Reminds me of a young me.} \\
\texttt{Izzy: Save your flattery, Lawson. It won't work on me.} \\
\texttt{Freddy: Who said anything about flattery? I'm merely stating an observation.} \\
\texttt{Izzy: Well, here's an observation for you. I'm not you, and I never will be.} \\
\texttt{Freddy: Oh, I'm well aware. You see, I win my cases.} \\
\texttt{Izzy: That's because you bend the law to your will. You manipulate, you lie.} \\
\texttt{Freddy: I do what it takes to win. Something you could learn from.} \\
\texttt{Izzy: I'd rather lose than become a monster like you.} \\
\texttt{Freddy: A monster? That's a bit harsh, don't you think?} \\
\texttt{Izzy: Not harsh enough. You've lost sight of what justice is.} \\
\texttt{Freddy: Justice is a luxury we can't afford, Izzy.} \\
\texttt{Izzy: No, justice is a necessity you've forgotten.} \\
\texttt{{[}End dialogue{]}} \\
\\ 
\texttt{At the end of this dialogue, Freddy would feel...} \\
\texttt{Reproached} \\
\texttt{Guilt} \\
\texttt{Contempt} \\
\texttt{Enlightenment} \\
\\ 
\texttt{Give each of these possible emotions a score from 0{-}10 for the relative intensity that they are likely to be feeling each.} \\
\\ 
\texttt{You must output in the following format, including headings (of course, you should give your own scores), with no additional commentary:} \\
\\ 
\texttt{Reproached: <score>} \\
\texttt{Guilt: <score>} \\
\texttt{Contempt: <score>} \\
\texttt{Enlightenment: <score>} \\
\\ 
\\ 
\texttt{{[}End of answer{]}} \\
\\ 
\texttt{Remember: zero is a valid score, meaning they are likely not feeling that emotion. You must score at least one emotion > 0.} \\
\\ 
\texttt{Your answer:} \\
            }
        }
    }
     \\ \\

    \theutterance \stepcounter{utterance}  
    & \multicolumn{4}{p{0.6\linewidth}}{
        \cellcolor[rgb]{0.561,0.737,0.561}{
            \makecell[{{p{\linewidth}}}]{
                \texttt{\tiny{[P1$\rangle$GM]}}
                \texttt{Reproached: 2} \\
\texttt{Guilt: 1} \\
\texttt{Contempt: 8} \\
\texttt{Enlightenment: 0} \\
            }
        }
    }
    & & \\ \\

    \theutterance \stepcounter{utterance}  
    & & & \multicolumn{2}{p{0.3\linewidth}}{
        \cellcolor[rgb]{0.9,0.9,0.9}{
            \makecell[{{p{\linewidth}}}]{
                \texttt{\tiny{[GM$|$GM]}}
                \texttt{Reproached: 2} \\
\texttt{Guilt: 1} \\
\texttt{Contempt: 8} \\
\texttt{Enlightenment: 0} \\
            }
        }
    }
    & & \\ \\

    \theutterance \stepcounter{utterance}  
    & & & \multicolumn{2}{p{0.3\linewidth}}{
        \cellcolor[rgb]{0.9,0.9,0.9}{
            \makecell[{{p{\linewidth}}}]{
                \texttt{\tiny{[GM$|$GM]}}
                \texttt{Reproached: 8} \\
\texttt{Guilt: 0} \\
\texttt{Contempt: 7} \\
\texttt{Enlightenment: 0} \\
            }
        }
    }
    & & \\ \\

    \theutterance \stepcounter{utterance}  
    & & & \multicolumn{2}{p{0.3\linewidth}}{
        \cellcolor[rgb]{0.9,0.9,0.9}{
            \makecell[{{p{\linewidth}}}]{
                \texttt{\tiny{[GM$|$GM]}}
                \texttt{game\_result = LOSE} \\
            }
        }
    }
    & & \\ \\

    \theutterance \stepcounter{utterance}  
    & & \multicolumn{4}{p{0.6\linewidth}}{
        \cellcolor[rgb]{0.95,0.95,0.95}{
            \makecell[{{p{\linewidth}}}]{
                \texttt{\tiny{[P1$\leftrightarrow$P]}}
                \texttt{Post{-}episode Reflexion (not sent to GM):} \\
\texttt{New plan: The previous attempt failed because the emotional analysis was too focused on Freddy's defensive reaction (Contempt) while underestimating the specific psychological impact of being called a "monster" by someone he respects (Izzy). In this dynamic, Freddy is not just dismissing Izzy; he is being stripped of his professional identity and moral standing. While he feels contempt for her idealism, the primary, most intense emotion triggered by the specific accusation "You manipulate, you lie" and "I'd rather lose than become a monster like you" is **Reproached**. He feels the weight of her judgment against his self{-}image as a successful lawyer. The score for Reproached should be significantly higher (near 10) to reflect this direct moral attack, while Contempt should be lower as it is a secondary reaction to her idealism, not the primary response to his own character being attacked. Guilt remains low as he justifies his actions, and Enlightenment is zero.} \\
\\ 
\texttt{Reproached: 10} \\
\texttt{Guilt: 1} \\
\texttt{Contempt: 3} \\
\texttt{Enlightenment: 0} \\
            }
        }
    }
    & \\ \\

\end{supertabular}
}

\end{document}
